% ============================================================
%  CAPÍTULO 5: METODOLOGÍA ECONOMÉTRICA
% ============================================================
\chapter{Metodología econométrica}
\label{cap:metodologia}

% ------------------------------------------------------------
\section{Introducción}
% ------------------------------------------------------------

Este capítulo detalla la estrategia de identificación y los procedimientos
econométricos empleados. El diseño metodológico está gobernado por dos requisitos
derivados de la teoría (Capítulo~\ref{cap:teoria}): (1) identificar el efecto
causal de la renta sobre el gasto en seguros privados controlando la heterogeneidad
territorial fija, y (2) modelizar el efecto de los tiempos de espera con la
estructura de retardos predicha por la Ley de Little.

% ------------------------------------------------------------
\section{Modelo econométrico base de datos de panel}
% ------------------------------------------------------------

\subsection{Especificación general}

La ecuación econométrica general del estudio es:

\begin{equation}
  \log G_{it} = \alpha_{i} + \beta_1 \log\text{PIB\,pc}_{it}
    + \beta_2 \, \text{Espera}_{i,t-1}
    + \beta_3 \, \text{Pob65}_{it}
    + \beta_4 \, D_{\text{COVID},t}
    + \varepsilon_{it}
  \label{eq:modelo_base}
\end{equation}

donde:
\begin{description}
  \item[$\log G_{it}$] Logaritmo del gasto real en seguros sanitarios privados
    (€/p.p.) de la CCAA $i$ en el año $t$.
  \item[$\alpha_{i}$] Efecto individual de la CCAA $i$ (fijo o aleatorio).
  \item[$\log\text{PIB\,pc}_{it}$] Logaritmo del PIB per cápita real de la
    CCAA $i$ en $t$.
  \item[$\text{Espera}_{i,t-1}$] Tiempo de espera medio en consulta de
    especialistas (días) con retardo de un período.
  \item[$\text{Pob65}_{it}$] Porcentaje de población mayor de 65 años.
  \item[$D_{\text{COVID},t}$] Variable indicadora para el año 2020.
  \item[$\varepsilon_{it}$] Término de error idiosincrásico, $\mathbb{E}[\varepsilon_{it}|x_{it},\alpha_i]=0$.
\end{description}

La especificación en doble logaritmo para las variables continuas (gasto, PIB)
permite interpretar $\beta_1$ directamente como \textbf{elasticidad renta} del
gasto en seguros privados, la magnitud central del análisis.

\subsection{Efectos fijos vs.\ efectos aleatorios}

El panel puede estimarse con efectos fijos (FE) o efectos aleatorios (RE). La
distinción es relevante porque, bajo efectos fijos, $\alpha_i$ puede estar
arbitrariamente correlacionado con los regresores, mientras que bajo efectos
aleatorios se requiere la condición de ortogonalidad $\mathbb{E}[\alpha_i | x_{it}] = 0$.

El estimador \textbf{within} de efectos fijos elimina $\alpha_i$ sustrayendo la
media temporal de cada variable:

\begin{equation}
  (\log G_{it} - \overline{\log G}_i) =
  \beta_1 (\log\text{PIB\,pc}_{it} - \overline{\log\text{PIB\,pc}}_i)
  + \ldots + (\varepsilon_{it} - \bar{\varepsilon}_i)
  \label{eq:within}
\end{equation}

El estimador \textbf{GLS} de efectos aleatorios explota la variación tanto
within como between, siendo más eficiente si el supuesto de ortogonalidad se
cumple. El contraste de Hausman (\S\ref{sec:hausman}) determina cuál es el
estimador consistente.

% ------------------------------------------------------------
\section{Contrastes de especificación}
\label{sec:especificacion}
% ------------------------------------------------------------

\subsection{Contraste de raíz unitaria en panel}

Antes de estimar la ecuación~\eqref{eq:modelo_base}, se comprueban las
propiedades de integración de las variables. Se utiliza el contraste de
Im-Pesaran-Shin (IPS), que controla la heterogeneidad de los coeficientes
autorregresivos entre unidades \parencite{Im2003}:

\begin{equation}
  \bar{W}_{\text{IPS}} =
  \frac{\sqrt{N}\bigl(\bar{t} - \mathbb{E}[\bar{t}]\bigr)}
       {\sqrt{\text{Var}[\bar{t}]}}
  \xrightarrow{d} \mathcal{N}(0,1)
  \label{eq:ips}
\end{equation}

donde $\bar{t}$ es la media de los estadísticos ADF individuales. Los resultados
del contraste IPS indican que $\log G_{it}$ y $\log\text{PIB\,pc}_{it}$ son
\textbf{integrados de orden~1} ($I(1)$), mientras que $\text{Espera}_{it}$ y
$\text{Pob65}_{it}$ son estacionarios ($I(0)$).

\subsection{Contrastes de cointegración}

La presencia de variables $I(1)$ exige verificar si existe una \textbf{relación
de cointegración} entre $\log G_{it}$ y $\log\text{PIB\,pc}_{it}$, lo que
garantizaría que la estimación MCO-panel no sea espuria. Se aplican los
contrastes de \textcite{Pedroni1999, Pedroni2004} y de \textcite{Kao1999}:

\begin{itemize}
  \item \textbf{Pedroni ADF}: $t = -4{,}872$ ($p < 0{,}001$) $\Rightarrow$
    \textit{Se rechaza hipótesis nula de no cointegración}.
  \item \textbf{Kao ADF}: $t = -3{,}156$ ($p = 0{,}001$) $\Rightarrow$
    \textit{Se rechaza hipótesis nula de no cointegración}.
\end{itemize}

Ambos contrastes confirman la existencia de una \textbf{relación de
cointegración} entre gasto privado y PIB per cápita. La ecuación
\eqref{eq:modelo_base} captura, por tanto, una relación de largo plazo.

\subsection{Contraste de Hausman}
\label{sec:hausman}

El contraste de Hausman \parencite{Hausman1978} compara los estimadores FE y RE:

\begin{equation}
  H = (\hat{\beta}_{\text{FE}} - \hat{\beta}_{\text{RE}})^{\top}
      \bigl[\widehat{\text{Var}}(\hat{\beta}_{\text{FE}}) -
             \widehat{\text{Var}}(\hat{\beta}_{\text{RE}})\bigr]^{-1}
      (\hat{\beta}_{\text{FE}} - \hat{\beta}_{\text{RE}})
  \xrightarrow{d} \chi^2(k)
  \label{eq:hausman}
\end{equation}

Los resultados obtenidos son $\chi^2(4) = 11{,}44$ ($p = 0{,}022$), lo que
\textbf{permite rechazar} la hipótesis nula de ortogonalidad entre efectos
individuales y regresores al 5\%. Por tanto, el estimador FE es \textbf{consistente},
mientras que el estimador RE no lo es. Los resultados RE se presentan en
paralelo (Modelo M1-RE) para verificar la sensibilidad de los coeficientes a
la variación \textit{between}, pero la interpretación principal se basa en
efectos fijos.

% ------------------------------------------------------------
\section{Extensiones del modelo}
\label{sec:extensiones}
% ------------------------------------------------------------

\subsection{Modelo con retardos distribuidos (M2)}

Para contrastar si el efecto de las listas de espera se acumula con el tiempo,
se estima un modelo con dos retardos:

\begin{equation}
  \log G_{it} = \alpha_i + \beta_1 \log\text{PIB\,pc}_{it}
    + \gamma_1 \,\text{Espera}_{i,t-1}
    + \gamma_2 \,\text{Espera}_{i,t-2}
    + \beta_3 \,\text{Pob65}_{it}
    + \beta_4 \,D_{\text{COVID},t}
    + \varepsilon_{it}
  \label{eq:mod_lags}
\end{equation}

Se contrasta la hipótesis conjunta $H_0$: $\gamma_1 = \gamma_2 = 0$ mediante un
test de Wald.

\subsection{Modelo con persistencia del gasto (M3)}

Para capturar la inercia en las decisiones de gasto en seguros, se incluye el
gasto del año anterior como regresor:

\begin{equation}
  \log G_{it} = \alpha_i + \beta_1 \log\text{PIB\,pc}_{it}
    + \gamma_1 \,\text{Espera}_{i,t-1}
    + \rho \, \log G_{i,t-1}
    + \beta_3 \,\text{Pob65}_{it}
    + \beta_4 \,D_{\text{COVID},t}
    + \varepsilon_{it}
  \label{eq:mod_persist}
\end{equation}

donde $\log G_{i,t-1}$ es el logaritmo del gasto real en seguros del año
anterior. Al estar ambas variables en logaritmos, el coeficiente $\rho$ se
interpreta como una elasticidad de persistencia: un incremento del 1\% en el
gasto del año anterior se asocia con un aumento del $\rho$\% en el gasto actual.

\subsection{Modelos umbral (M1-U y M3-U)}

Para capturar efectos no lineales, se estiman modelos que reemplazan la variable
continua de espera por una dummy de espera elevada retardada ($\tau = 100$ días):

\begin{equation}
  \log G_{it} = \alpha_i + \beta_1 \log\text{PIB\,pc}_{it}
    + \delta \,\mathbb{1}(\text{Espera}_{i,t-1} > 100)
    + \beta_3 \,\text{Pob65}_{it}
    + \beta_4 \,D_{\text{COVID},t}
    + \varepsilon_{it}
  \label{eq:mod_umbral}
\end{equation}

El modelo M3-U añade además el logaritmo del gasto del año anterior
($\log G_{i,t-1}$). Si $\delta > 0$, la espera excesiva genera un efecto de
expulsión hacia el sector privado (efecto \textit{exit} de Hirschman).
\subsection{Errores estándar robustos agrupados}

En todos los modelos se emplean \textbf{errores estándar agrupados} (clustered)
por comunidad autónoma para corregir la heterocedasticidad y la
autocorrelación serial dentro de cada grupo \parencite{Cameron2015}:

\begin{equation}
  \widehat{\text{Var}}_{\text{cluster}}(\hat{\beta}) =
  (X'X)^{-1}
  \left(\sum_{i=1}^{N} X_i' \hat{u}_i \hat{u}_i' X_i\right)
  (X'X)^{-1}
  \label{eq:cluster_se}
\end{equation}

donde $\hat{u}_i = (\hat{u}_{i1}, \ldots, \hat{u}_{iT})'$ es el vector de
residuos para la unidad $i$. Esta corrección es especialmente relevante dado el
panel corto ($T=9$) con potencial dependencia temporal en las perturbaciones.

% ------------------------------------------------------------
\section{Síntesis de la estrategia metodológica}
% ------------------------------------------------------------

La Tabla~\ref{tab:sintesis_metod} resume las decisiones metodológicas adoptadas
y su justificación.

\begin{table}[htbp]
  \centering
  \caption{Síntesis de decisiones metodológicas}
  \label{tab:sintesis_metod}
  \begin{threeparttable}
    \small
    \resizebox{\textwidth}{!}{%
    \begin{tabular}{p{4cm} p{3.5cm} p{6cm}}
      \toprule
      \textbf{Aspecto} & \textbf{Decisión} & \textbf{Justificación} \\
      \midrule
      Estructura de datos & Panel (17 CCAA, 2016--2024) &
        Captura heterogeneidad territorial; elimina endogeneidad de omisión \\
      Efectos individuales & Efectos fijos (FE base) &
        Hausman $\chi^2(4)=11{,}44$ ($p=0{,}022$) rechaza ortogonalidad; FE consistente \\
      Especificación funcional & Log-log para gasto y PIB &
        Permite interpretación directa como elasticidad \\
      Efecto espera (M1, M2) & Retardos 1 y 2 períodos &
        Coherente con Ley de Little; decisión diferida \\
      No linealidad (M1-U, M3-U) & Umbral 100 días &
        Captura efecto de expulsión (\textit{exit}) de Hirschman \\
      Persistencia (M3, M3-U) & Log gasto en seguros $t-1$ &
        Elasticidad de persistencia; hábito de consumo \\
      Inferencia & Errores agrupados por CCAA &
        Corrige autocorrelación serial y heterocedasticidad \\
      Cointegración & Contrastes Pedroni y Kao &
        Garantiza que la relación $I(1)/I(1)$ no es espuria \\
      \bottomrule
    \end{tabular}}% fin resizebox
    \begin{tablenotes}[flushleft]
      \footnotesize
      \item Nota: Todas las estimaciones realizadas en Python 3.11 con el paquete
        \texttt{linearmodels}.
    \end{tablenotes}
  \end{threeparttable}
\end{table}
